% =====================================================================
% Wave-19 Z1/20 -- Borcea + Voisin voices.
% $N=5$ non-CHL effective positive-half geometry on the Borcea--Voisin
% threefold $X_{5}^{\mathrm{BV}} = (S_{5}\times E_{5})/(\iota_{S}\times\iota_{E})$,
% with metaplectic fourth-root character
% $(Z^{X_{5}^{\mathrm{BV}}})^{-1/4}=\Delta_{1/2}^{(5)}$.
% CG voice. Subscripted kappa only. \providecommand macros only.
% Primary sources: Borcea 1997 (K3 with involution + mirror pairs);
%   Voisin 1993 (miroirs et involutions); Borcherds 1998 Invent. Math. 132
%   Thm 10.1, Thm 13.3 (metaplectic singular-theta lift + weight formula);
%   Gritsenko-Clery 2008 arXiv:0812.3962 Thm 1.1 row N=5; Nikulin 1980
%   Trudy Mosk. Mat. Obs. 38 Thm 4.5 (admissible orders); Nikulin 1979
%   Prop 4.2 (fixed-locus genera for antisymplectic involution);
%   Mukai 1988 Invent. Math. 94 Thm 0.6 + Table 1 (Niemeier attachments,
%   #Fix(g_5)=4); Mongardi-Tari-Wandel 2018 arXiv:1610.06216 Remark 2.4
%   (no Kummer-HK for order 5); Maulik-Pandharipande-Thomas 2010
%   (reduced virtual class on K3 x E); Oberdieck-Pixton 2018
%   arXiv:1808.03524 Conj. A (DT zeta on K3 x E); Gritsenko 1999
%   St. Petersburg Math. J. 11 Thm 1.2 (c_5(0)=1, psi_{1/2,5} construction);
%   Lorgat 2020 Sec 6 (Borcherds lift template at N=1).
% Author: Raeez Lorgat.
% =====================================================================

\providecommand{\ZZ}{\mathbb{Z}}
\providecommand{\QQ}{\mathbb{Q}}
\providecommand{\CC}{\mathbb{C}}
\providecommand{\HH}{\mathbb{H}}
\providecommand{\kBKM}{\kappa_{\mathrm{BKM}}}
\providecommand{\kch}{\kappa_{\mathrm{ch}}}
\providecommand{\kcat}{\kappa_{\mathrm{cat}}}
\providecommand{\kfib}{\kappa_{\mathrm{fiber}}}
\providecommand{\Aut}{\mathrm{Aut}}
\providecommand{\Auts}{\mathrm{Aut}_{s}}
\providecommand{\Fix}{\mathrm{Fix}}
\providecommand{\Mp}{\mathrm{Mp}}
\providecommand{\Sp}{\mathrm{Sp}}
\providecommand{\wt}{\mathrm{wt}}
\providecommand{\BV}{\mathrm{BV}}
\providecommand{\BPS}{\mathrm{BPS}}
\providecommand{\DT}{\mathrm{DT}}
\providecommand{\red}{\mathrm{red}}
\providecommand{\Delfnm}[2]{\Delta_{#1}^{(#2)}}
\providecommand{\gBPS}{\mathfrak{g}_{\BPS}}
\providecommand{\ClaimStatusTheorem}{\textsc{[Thm]}}
\providecommand{\ClaimStatusConjectured}{\textsc{[Conj]}}
\providecommand{\ClaimStatusHeuristic}{\textsc{[Heur]}}

\section*{Wave-19 Z1/20 --- $N=5$ Borcea--Voisin host for $\Delfnm{1/2}{5}$}
\label{wn:sec:wave19-z1-N5-borcea-voisin}

$\varphi(5)=4\nmid 2$: no elliptic curve admits $\ZZ[\zeta_{5}]
\hookrightarrow\mathrm{End}(E)$, so no CHL diagonal $\ZZ/5$ on
$K3\times E$. The paramodular form $\Delfnm{1/2}{5}\in
M_{1/2}(K(5),\chi_{5})$ of weight $1/2$ exists (Gritsenko--Cl\'ery
2008 Thm~1.1 row~$N=5$). Its chiral host is the Borcea--Voisin
threefold $X_{5}^{\BV}$, with $\ZZ/5$ acting on the K3 factor and
the elliptic factor absorbing an order-two involution; the identity
$(Z^{X_{5}^{\BV}})^{-1/4}=\Delfnm{1/2}{5}$ is forced by the square
of weight $1/2$ landing on weight $1$, the fourth power on weight
$2$ on $\Sp_{4}(\ZZ)$.

\subsection*{Cycle 1 --- Borcea--Voisin threefold} \ClaimStatusTheorem

$S_{5}$ is a K3 surface carrying a symplectic order-$5$ automorphism
$g_{5}\in\Auts(S_{5})$; Nikulin 1980 Thm~4.5 lists admissible cyclic
symplectic orders $N\in\{1,\dots,8\}$, and at $N=5$ the fixed count
is $\#\,\Fix(g_{5})=4$ with Mukai coinvariant Niemeier anchor
$4A_{4}$ of rank $16$ (Mukai 1988 Thm~0.6, Table~1 row~$\ZZ/5$).
$S_{5}$ further carries an anti-symplectic Nikulin involution
$\iota_{S}$ commuting with $g_{5}$; its fixed locus is a disjoint
union of smooth curves $\Fix(\iota_{S})=\sqcup_{j=1}^{k}C_{j}$ whose
numerics are governed by the Nikulin triple $(r,a,\delta)\in\ZZ^{3}$
via
\[
k\,+\,\textstyle\sum_{j}g(C_{j})\;=\;11-r,\qquad
k\,-\,\textstyle\sum_{j}g(C_{j})\;=\;a
\]
(Nikulin 1979 Prop.~4.2, Voisin 1993 \S 1.2). Let $E_{5}$ be an
elliptic curve with $\iota_{E}(z)=-z$. The Borcea--Voisin threefold
\[
X_{5}^{\BV}\;:=\;\bigl(S_{5}\times E_{5}\bigr)\big/(\iota_{S}\times\iota_{E}),
\]
crepantly resolved $\pi:\widehat{X}_{5}^{\BV}\to X_{5}^{\BV}$ along
$\Fix(\iota_{S})\times\Fix(\iota_{E})$ via $\mathbb{P}^{1}$-bundles over
the fixed curves (Borcea 1997 Thm~1; Voisin 1993 Thm~2.1), is smooth
Calabi--Yau of dimension $3$ with
$\chi_{\mathrm{top}}(\widehat{X}_{5}^{\BV})=2(r-10)(11-r)-8a$
(Voisin 1993 Cor.~2.2). The symplectic $\ZZ/5=\langle g_{5}\rangle$
descends to $\widehat{X}_{5}^{\BV}$, acting freely on a Zariski-open
locus with fixed locus the finite set
$\Fix(g_{5})\cap\Fix(\iota_{S})\subset S_{5}$ paired with
$\Fix(\iota_{E})=\{0,P_{2}\}\subset E_{5}$.

\subsection*{Cycle 2 --- effective positive-half moduli} \ClaimStatusConjectured

The moduli stack $\mathcal{M}_{\beta,n}(\widehat{X}_{5}^{\BV})$ of
$(g_{5},\iota_{S}\times\iota_{E})$-equivariant stable sheaves carries
a reduced virtual class
$[\mathcal{M}]^{\red}\in A_{0}^{\ZZ/5}(\mathcal{M})$: both $g_{5}$
and $\iota_{S}\times\iota_{E}$ preserve the product holomorphic
$3$-form up to fifth and sign roots of unity, yielding a Maulik--Toda
cosection on $\mathcal{M}$ and a $(\ZZ/5\times\ZZ/2)$-equivariant
reduction after Maulik--Pandharipande--Thomas 2010 Thm~1 (class
construction on $K3\times E$) extended to Borcea--Voisin twist
(cf.\ Bryan--Steinberg 2014 Thm~1.1 for crepant symplectic orbifolds).
The positive-half generating series
\[
\mathsf{Z}^{X_{5}^{\BV}}_{\beta}(p,q,y)\;:=\;
\sum_{n\geq 0}\,p^{n}\,q^{\beta^{2}/2}\,y^{(\beta,f)}\,
\int_{[\mathcal{M}_{\beta,n}]^{\red}}\!\mathbf{1}
\]
is conjectured (Oberdieck--Pixton 2018 Conj.~A pattern) to be a
weight-$2$ meromorphic Siegel form on $K(5)$ with divisor on the
Humbert surface $\{z=0\}\cup\{z=\tau/5\}\cup\ldots$, satisfying
\[
\mathsf{Z}^{X_{5}^{\BV}}_{\beta=L}(p,q,y)
\;=\; -C_{5}\cdot\bigl(\Delfnm{1/2}{5}(p,q,y)\bigr)^{-4},
\qquad C_{5}\in\QQ^{\times}.
\]
The exponent $-4$ is forced by Cycle~4.

\subsection*{Cycle 3 --- metaplectic fourth-root character} \ClaimStatusTheorem

$\Delfnm{1/2}{5}$ has weight $1/2$ under the character $\chi_{5}$ of
the metaplectic cover $\Mp_{4}(\ZZ)\twoheadrightarrow\Sp_{4}(\ZZ)$
restricted to $K(5)$ (Borcherds 1998 Thm~10.1; GC~2008 Thm~1.1 row
$N=5$). The square $(\Delfnm{1/2}{5})^{2}\in M_{1}(K(5),
\chi_{5}^{\otimes 2})$ has trivial character on the double cover;
the fourth power $(\Delfnm{1/2}{5})^{4}\in M_{2}(K(5))$ descends to
a genuine paramodular cusp form of weight $2$. The singular-theta
weight formula (Borcherds 1998 Thm~13.3) gives
$\wt(\Delfnm{1/2}{5})=c_{5}(0)/2=1/2$ from input Jacobi form
$\psi_{1/2,5}=\vartheta_{1}\,\eta^{-3}f_{5}$ with $f_{5}$ the weight-$2$
cusp generator of $S_{2}(\Gamma_{0}(5))$ (Gritsenko 1999 Thm~1.2,
\S 3), whose $(q^{0}\zeta^{0})$-coefficient is $c_{5}(0)=1$. Hence
\[
\bigl(\mathsf{Z}^{X_{5}^{\BV}}\bigr)^{-1/4}\;=\;\Delfnm{1/2}{5}
\]
is the unique root landing on a weight-$1$ form after squaring, a
weight-$2$ paramodular form after fourth-powering: a metaplectic
square-root for $\Mp_{4}\to\Sp_{4}$, composed with the further
$\chi_{5}$-doubling of primitive order $4$ on $K(5)$.

\subsection*{Cycle 4 --- universal Borcherds identity} \ClaimStatusTheorem

The extension $\kBKM(\Phi_{N})=c_{N}(0)/2$ to $N=5$ is a
Borcherds-weight theorem: the input $\psi_{1/2,5}\in
J_{1/2,5}^{w,\chi_{5}}$ has $c_{5}(0)=1$ and Borcherds' formula
reads $\wt(\Delfnm{1/2}{5})=1/2$ directly (Wave-16 V4 Cycle~1).
The BPS-algebra interpretation is new at $N=5$: the Gritsenko--Nikulin
1998 Part~II denominator framework covers $N\in\{1,2,3,4,6\}$, and
the Borcea--Voisin host supplies the chiral-side realisation at
$N=5$. The Borcherds multiplicative lift
\[
\Delfnm{1/2}{5}(Z)\;=\;\prod_{(m,n,\ell)>0}\bigl(1-e(m\tau+n\sigma+\ell z)\bigr)^{c_{5}(mn,\ell)}
\]
converges absolutely on $\{(\Im\tau)(\Im\sigma)-(\Im z)^{2}>\tfrac14\}
\subset\HH_{2}$; the divisor $\{z=0\}$ has multiplicity
$c_{5}(0,1)=-2$, consistent with the pair-wise Weyl-chamber count
for $\#\,\Fix(g_{5})=4$ \ClaimStatusHeuristic{} (no primary
derivation of the $-2$ to fixed-count bridge).

\subsection*{Cycle 5 --- $\gBPS^{(5)}$ and automorphic identification} \ClaimStatusConjectured

The BPS algebra $\gBPS^{(5)}:=\gBPS(\widehat{X}_{5}^{\BV})$ is the
cohomological Hall algebra of $g_{5}$-equivariant stable sheaves on
$\widehat{X}_{5}^{\BV}$ with reduced virtual Euler class
(Kontsevich--Soibelman 2010 \S 6 adapted from $K3\times E$). Its
generating character
\[
\chi\bigl(\gBPS^{(5)}\bigr)(p,q,y)\;=\;
-\log\mathsf{Z}^{X_{5}^{\BV}}(p,q,y)
\;=\;4\log\Delfnm{1/2}{5}(p,q,y)
\]
matches the root-multiplicity expansion of the metaplectic GKM
superalgebra $\mathfrak{g}_{\Delfnm{1/2}{5}}$ attached to
$\Delfnm{1/2}{5}$ by the Borcherds--Gritsenko--Nikulin correspondence
in its metaplectic extension (Borcherds 1998 \S 14; rigorous at
$N=5$ pending metaplectic BKM denominator identity). The
identification
\[
\gBPS^{(5)}\;\cong\;\mathfrak{g}_{\Delfnm{1/2}{5}}
\]
is the $N=5$ analogue of Oberdieck's twisted Mathieu theorem on
$K3\times E$; its proof reduces to (i) reduced virtual class on
$\mathcal{M}_{\beta,n}(\widehat{X}_{5}^{\BV})$ with
$(g_{5},\iota)$-equivariance (MPT 2010 + Borcea--Voisin twist);
(ii) $K(5)$-modularity with character $\chi_{5}$ (Oberdieck
modularity extended to non-CHL by fourth-root convention); (iii)
Fourier matching against Cycle~4.

\subsection*{Kappa data}

$\kcat(\widehat{X}_{5}^{\BV})=\chi(\mathcal{O})=1-0+0-1=0$: the
Hodge supertrace on a smooth CY$_{3}$, equivalently $\kch=0$ at
$d=3$ where $\kch=\chi(\mathcal{O})$ on compact CY$_{d}$.
$\kfib$ records the orbifold correction from the exceptional
$\mathbb{P}^{1}$-bundles $\pi^{-1}(C_{j}\times\Fix(\iota_{E}))$:
each fixed curve of genus $g_{j}$ contributes $(1-g_{j})$, summing
to $\kfib=\sum_{j}(1-g_{j})=k-(11-r-k)=2k+r-11$ via Nikulin's
$\sum g_{j}+k=11-r$. $\kBKM(\Delfnm{1/2}{5})=1/2$ lands outside
the $\{\kch+\kfib\}$ decomposition of the CHL slice: at $N=5$ the
identity $\kBKM=c_{5}(0)/2$ is a singular-theta weight statement
(metaplectic scope), whereas $\{\kch,\kcat,\kfib\}$ remain classical
Hodge data of $\widehat{X}_{5}^{\BV}$.

\subsection*{Scope}

The threefold $X_{5}^{\BV}$ and its crepant resolution are theorems
(Borcea 1997; Voisin 1993). The Borcherds weight $\wt(\Delfnm{1/2}{5})
=c_{5}(0)/2=1/2$ is a theorem (Borcherds 1998 Thm~13.3; Gritsenko
1999). The fourth-root identity
$(\mathsf{Z}^{X_{5}^{\BV}})^{-1/4}=\Delfnm{1/2}{5}$ is conjectural,
resting on Oberdieck--Pixton-type modularity of the reduced DT zeta on a
non-CHL host. The identification $\gBPS^{(5)}\cong
\mathfrak{g}_{\Delfnm{1/2}{5}}$ is conjectural pending a metaplectic
BKM denominator identity at level $5$.

\paragraph{Cite.} Borcea, \emph{K3 surfaces with involution and
mirror pairs of CY manifolds}, Essays on Mirror Manifolds II (1997)
717--743, Thm~1; Voisin, \emph{Miroirs et involutions sur les surfaces
K3}, Ast\'erisque~218 (1993) 273--323, Thm~2.1, Cor.~2.2;
Borcherds, Invent.\ Math.~\textbf{132} (1998) Thm~10.1, Thm~13.3;
Gritsenko--Cl\'ery arXiv:0812.3962 (2008) Thm~1.1 row~$N=5$;
Nikulin, Trudy Mosk.\ Mat.\ Obs.~\textbf{38} (1980) Thm~4.5;
Nikulin, \emph{Izv.\ Akad.\ Nauk}~\textbf{43} (1979) Prop.~4.2;
Mukai, Invent.\ Math.~\textbf{94} (1988) Thm~0.6, Table~1;
Mongardi--Tari--Wandel arXiv:1610.06216 (2018) Remark~2.4;
Maulik--Pandharipande--Thomas, J.~Topol.~\textbf{3} (2010) Thm~1;
Oberdieck--Pixton arXiv:1808.03524 (2018) Conj.~A;
Gritsenko, St.\ Petersburg Math.\ J.~\textbf{11} (1999) Thm~1.2;
Kontsevich--Soibelman arXiv:1006.2706 (2010) \S 6;
Lorgat 2020 \S 6.
